\chapter{Competition Participation}

\section{Dates and Venue}

The CAuDri-Challenge \yearofchallenge{} will take place on \dateofchallenge{} at the DHBW Stuttgart Fakultät Technik (Lerchenstraße 1, 70174 Stuttgart).
The exact schedule will be published on our website.


\section{Prerequisites for Attending}

Only students and teams fulfilling the following conditions are allowed to participate in the CAuDri-Challenge.

\subsection{Status of Enrollment}

Every participant must either be currently enrolled in a Bachelor’s, Master’s or a comparable degree program or the respective degree must not have been obtained more than six months before the competition.
There is no restriction concerning the subject of study.
Research staff and PhD students may not participate actively in conceptualization or development of the vehicle.
They may not participate actively in the competition (cf. Section \ref{dev_know_how}).

While no proof of enrollment is required for registration, we reserve the right to request a proof of enrollment at any time.
Violations of this rule will result in the disqualification of the respective team.

\subsection{Minimum Age}

There is no minimum age for participation. Underage participants must submit a non formal written declaration of consent signed by their legal guardian.
A template can be found on our \href{https://caudri-challenge.de/documents/Einverst%C3%A4ndniserkl%C3%A4rung_Minderj%C3%A4hrige.pdf}{website}.

\subsection{Number of Teams per Institution}

The number of teams per institution is not limited.
However, the development of the vehicles must be strictly separated. Software and hardware architectures of the respective teams must differ significantly.

\subsection{Publication Rights}

By registering, every team and every participant declares their agreement with the publication of image, video and audio recordings.
This also includes the recording of team presentations.
This agreement might be revoked until the day of the competition.

% \section{Limit of participating teams}

% The commission reserves the right to limit the number of participating teams if
% this becomes necessary due to organizational reasons. In case of such a limit
% additional qualifying tasks may be requested from the teams.